\documentclass[conference]{IEEEtran}
\IEEEoverridecommandlockouts
%Template version as of 6/27/2024

\usepackage{cite}
\usepackage{amsmath,amssymb,amsfonts}
\usepackage{algorithmic}
\usepackage{graphicx}
\usepackage{textcomp}
\usepackage{xcolor}
\usepackage[utf8]{inputenc}
\usepackage[T1]{fontenc}
\usepackage{booktabs}

\def\BibTeX{{\rm B\kern-.05em{\sc i\kern-.025em b}\kern-.08em
    T\kern-.1667em\lower.7ex\hbox{E}\kern-.125emX}}

\begin{document}

\title{Detección Automática de Fracturas Óseas en Radiografías Mediante Modelos de Visión por Computadora Basados en Deep Learning}

\author{
\IEEEauthorblockN{Lucia T. Capon Paul}
\IEEEauthorblockA{\textit{CEIA -- FIUBA} \\
Buenos Aires, Argentina}
\and
\IEEEauthorblockN{Cesar Orellana}
\IEEEauthorblockA{\textit{CEIA -- FIUBA} \\
Buenos Aires, Argentina}
\and
\IEEEauthorblockN{Leandro Britez}
\IEEEauthorblockA{\textit{CEIA -- FIUBA} \\
Buenos Aires, Argentina}
}

\maketitle

% ÔöÇÔöÇÔöÇ ABSTRACT ÔöÇÔöÇÔöÇÔöÇÔöÇÔöÇÔöÇÔöÇÔöÇÔöÇÔöÇÔöÇÔöÇÔöÇÔöÇÔöÇÔöÇÔöÇÔöÇÔöÇÔöÇÔöÇÔöÇÔöÇÔöÇÔöÇÔöÇÔöÇÔöÇÔöÇÔöÇÔöÇÔöÇÔöÇÔöÇÔöÇÔöÇÔöÇÔöÇÔöÇÔöÇÔöÇÔöÇÔöÇÔöÇÔöÇÔöÇÔöÇÔöÇÔöÇÔöÇÔöÇÔöÇÔöÇÔöÇÔöÇÔöÇÔöÇÔöÇÔöÇÔöÇÔöÇÔöÇÔöÇ
\begin{abstract}
% TODO: completar tras finalizar los experimentos (~150 palabras)
% Incluir: problema, enfoque, modelos comparados (YOLOv11 vs RT-DETR),
%          dataset usado, m├®tricas clave obtenidas, y principal conclusi├│n.
La detección automatizada de fracturas óseas en imágenes radiográficas representa
un desafío clínico de alto impacto. En este trabajo se presenta el desarrollo de un
prototipo basado en modelos de detecci├│n de objetos preentrenados para la
identificación y localización de fracturas en radiografías.
Se comparan dos arquitecturas: YOLOv11, representativa del paradigma
anchor-free de una etapa, y RT-DETR, un detector basado en transformers de
tiempo real. Los modelos son ajustados (\textit{fine-tuned}) sobre el dataset
p├║blico \textit{Bone Fracture Detection} de Roboflow Universe.
La evaluaci├│n se realiza en t├®rminos de mAP@0.5, mAP@0.5:0.95, F1-score y
latencia de inferencia. [COMPLETAR CON RESULTADOS].
\end{abstract}

\begin{IEEEkeywords}
detecci├│n de objetos, fracturas ├│seas, YOLOv11, RT-DETR, deep learning,
visión por computadora, radiografías.
\end{IEEEkeywords}

% ÔöÇÔöÇÔöÇ I. INTRODUCCI├ôN ÔöÇÔöÇÔöÇÔöÇÔöÇÔöÇÔöÇÔöÇÔöÇÔöÇÔöÇÔöÇÔöÇÔöÇÔöÇÔöÇÔöÇÔöÇÔöÇÔöÇÔöÇÔöÇÔöÇÔöÇÔöÇÔöÇÔöÇÔöÇÔöÇÔöÇÔöÇÔöÇÔöÇÔöÇÔöÇÔöÇÔöÇÔöÇÔöÇÔöÇÔöÇÔöÇÔöÇÔöÇÔöÇÔöÇÔöÇÔöÇÔöÇÔöÇÔöÇÔöÇÔöÇÔöÇÔöÇÔöÇÔöÇ
\section{Introducci├│n}

La interpretación de imágenes radiográficas es una tarea clínica fundamental en
el diagnóstico de traumatismos y patologías óseas. La identificación visual de
fracturas requiere experiencia especializada y puede verse afectada por factores
como la carga de trabajo y la fatiga del profesional \cite{b1}.

Los sistemas de visi├│n por computadora basados en \textit{deep learning} han
demostrado capacidad para asistir en tareas de detecci├│n y localizaci├│n de
patolog├¡as en im├ígenes m├®dicas \cite{b2}. En particular, los modelos de
detecci├│n de objetos permiten no solo determinar la presencia de una lesi├│n,
sino tambi├®n localizar su posici├│n mediante cuadros delimitadores
(\textit{bounding boxes}).

% TODO: agregar caso real de aplicaci├│n similar (no mayor a 8 a├▒os)

El objetivo de este trabajo es desarrollar y evaluar un prototipo funcional
capaz de detectar automáticamente fracturas óseas en radiografías, comparando
el desempe├▒o de dos arquitecturas modernas: YOLOv11 \cite{b3} y RT-DETR \cite{b4}.

% ÔöÇÔöÇÔöÇ II. METODOLOG├ìA ÔöÇÔöÇÔöÇÔöÇÔöÇÔöÇÔöÇÔöÇÔöÇÔöÇÔöÇÔöÇÔöÇÔöÇÔöÇÔöÇÔöÇÔöÇÔöÇÔöÇÔöÇÔöÇÔöÇÔöÇÔöÇÔöÇÔöÇÔöÇÔöÇÔöÇÔöÇÔöÇÔöÇÔöÇÔöÇÔöÇÔöÇÔöÇÔöÇÔöÇÔöÇÔöÇÔöÇÔöÇÔöÇÔöÇÔöÇÔöÇÔöÇÔöÇÔöÇÔöÇÔöÇÔöÇÔöÇÔöÇÔöÇ
\section{Metodología, Diseño y Desarrollo}

\subsection{Dataset}

Se utiliz├│ el dataset \textit{Bone Fracture Detection} disponible p├║blicamente
en Roboflow Universe \cite{b5}. El dataset contiene imágenes de radiografías con
anotaciones en formato YOLO (\textit{bounding boxes}). Las particiones utilizadas
son:

\begin{table}[htbp]
\caption{Distribuci├│n del dataset por split}
\begin{center}
\begin{tabular}{lcc}
\toprule
\textbf{Split} & \textbf{Imágenes} & \textbf{Anotaciones} \\
\midrule
Train & -- & -- \\  % TODO: completar con valores reales
Valid & -- & -- \\
Test  & -- & -- \\
\bottomrule
\end{tabular}
\end{center}
\end{table}

\subsection{Preprocesamiento y Data Augmentation}

Todas las imágenes fueron redimensionadas a $640 \times 640$ píxeles. Se
aplicaron las siguientes t├®cnicas de \textit{data augmentation} durante
el entrenamiento: variaciones de matiz y saturaci├│n (HSV), flip horizontal,
traslaci├│n aleatoria y escalado. Para RT-DETR se desactiv├│ el \textit{mosaic
augmentation} por ser incompatible con la arquitectura \textit{transformer}.

\subsection{Arquitecturas}

\textbf{YOLOv11} \cite{b3}: detector \textit{anchor-free} de una sola etapa.
Se utiliz├│ la variante \texttt{yolo11m.pt} preentrenada en COCO.

\textbf{RT-DETR} \cite{b4}: detector basado en transformers dise├▒ado para
inferencia en tiempo real. Se utiliz├│ la variante \texttt{rtdetr-l.pt}
preentrenada en COCO.

\subsection{Entrenamiento}

Ambos modelos fueron ajustados durante 100 ├®pocas con optimizador
\textit{AdamW}. Los hiperparámetros completos se documentan en el repositorio
de GitHub en los archivos \texttt{configs/yolov11.yaml} y
\texttt{configs/model2.yaml}.

% ÔöÇÔöÇÔöÇ III. RESULTADOS ÔöÇÔöÇÔöÇÔöÇÔöÇÔöÇÔöÇÔöÇÔöÇÔöÇÔöÇÔöÇÔöÇÔöÇÔöÇÔöÇÔöÇÔöÇÔöÇÔöÇÔöÇÔöÇÔöÇÔöÇÔöÇÔöÇÔöÇÔöÇÔöÇÔöÇÔöÇÔöÇÔöÇÔöÇÔöÇÔöÇÔöÇÔöÇÔöÇÔöÇÔöÇÔöÇÔöÇÔöÇÔöÇÔöÇÔöÇÔöÇÔöÇÔöÇÔöÇÔöÇÔöÇÔöÇÔöÇÔöÇÔöÇ
\section{Resultados Experimentales y Análisis}

% TODO: completar esta secci├│n tras correr los experimentos.

\begin{table}[htbp]
\caption{Comparaci├│n de m├®tricas en el conjunto de test}
\begin{center}
\begin{tabular}{lcccc}
\toprule
\textbf{Modelo} & \textbf{mAP@0.5} & \textbf{mAP@0.5:0.95} & \textbf{F1} & \textbf{Latencia (ms)} \\
\midrule
YOLOv11  & -- & -- & -- & -- \\
RT-DETR  & -- & -- & -- & -- \\
\bottomrule
\end{tabular}
\end{center}
\end{table}

% Agregar figuras de curvas PR, ejemplos de detección, análisis de errores.

% ÔöÇÔöÇÔöÇ IV. CONCLUSIONES ÔöÇÔöÇÔöÇÔöÇÔöÇÔöÇÔöÇÔöÇÔöÇÔöÇÔöÇÔöÇÔöÇÔöÇÔöÇÔöÇÔöÇÔöÇÔöÇÔöÇÔöÇÔöÇÔöÇÔöÇÔöÇÔöÇÔöÇÔöÇÔöÇÔöÇÔöÇÔöÇÔöÇÔöÇÔöÇÔöÇÔöÇÔöÇÔöÇÔöÇÔöÇÔöÇÔöÇÔöÇÔöÇÔöÇÔöÇÔöÇÔöÇÔöÇÔöÇÔöÇÔöÇÔöÇÔöÇÔöÇ
\section{Conclusiones}

% TODO: al menos 5 conclusiones (requerimiento de la cátedra)

\begin{enumerate}
    \item ...
    \item ...
    \item ...
    \item ...
    \item ...
\end{enumerate}

% ÔöÇÔöÇÔöÇ BIBLIOGRAF├ìA ÔöÇÔöÇÔöÇÔöÇÔöÇÔöÇÔöÇÔöÇÔöÇÔöÇÔöÇÔöÇÔöÇÔöÇÔöÇÔöÇÔöÇÔöÇÔöÇÔöÇÔöÇÔöÇÔöÇÔöÇÔöÇÔöÇÔöÇÔöÇÔöÇÔöÇÔöÇÔöÇÔöÇÔöÇÔöÇÔöÇÔöÇÔöÇÔöÇÔöÇÔöÇÔöÇÔöÇÔöÇÔöÇÔöÇÔöÇÔöÇÔöÇÔöÇÔöÇÔöÇÔöÇÔöÇÔöÇÔöÇÔöÇÔöÇÔöÇÔöÇ
\begin{thebibliography}{00}

\bibitem{b1}
% TODO: completar con paper de Papers/ relevante
Author, ``Title,'' Journal, vol. X, pp. Y--Z, year.

\bibitem{b2}
% TODO: completar con paper de Papers/
Author, ``Title,'' in Proc. Conference, year, pp. X--Y.

\bibitem{b3}
Ultralytics, ``YOLO11: State-of-the-Art Object Detection,'' 2024.
[Online]. Available: https://github.com/ultralytics/ultralytics

\bibitem{b4}
L. Leng et al., ``DETRs Beat YOLOs on Real-time Object Detection,''
arXiv:2304.08069, 2023.

\bibitem{b5}
Roboflow Universe, ``Bone Fracture Detection Dataset,'' 2024.
[Online]. Available: https://universe.roboflow.com/veda/bone-fracture-detection-daoon

\end{thebibliography}

\end{document}
